\section{Faster interchange of address chunks}
\label{sec:tape-interchange}

We use the bit network to interchange two equal address chunks on tape.
An edge matrix of rank $a$ will require $a$ smaller interchanges.
Because the sum of these ranks is less than the number of wires times
$m$, the resulting recurrence saves a power of the chunk width.

\subsection{A matrix decomposition adapted to tape order}

Lemma~\ref{lem:ordered-affine-streams} updates a later field using
earlier fields in linear time.  Lower triangular changes of coordinates
can be performed entirely in this direction. The required factorization
is the Bruhat factorization for arbitrary matrices, with coordinate order
reversed; see Grigor'ev~\cite[Appendix~A.2, Proposition~14(c)]{Grigoriev1982}.
We give its elementary elimination proof in this order.

\begin{lemma}[Lower triangular factorization]
\label{lem:lower-lower}
Every rational $m\times m$ matrix $A$ of rank $a$ has a factorization
\[
 A=E_1\Pi E_2,
\]
where $E_1,E_2$ are invertible lower triangular rational matrices and
$\Pi$ has exactly $a$ entries equal to one, with at most one in each row
and column, and all remaining entries zero.
\end{lemma}

\begin{proof}
If $A=0$, take $E_1=E_2=I$ and $\Pi=0$.  Otherwise choose the topmost
nonzero row and its rightmost nonzero entry.  Use that pivot column to
clear entries to its left by right multiplication by lower triangular
elementary matrices.  There are no nonzero entries to its right in the
pivot row.  Clear entries below the pivot by left multiplication by
lower triangular elementary matrices, and scale the pivot to one.
All rows above the chosen row were zero in the still active columns.

Delete the pivot row and column from the active index sets, retaining
their original orders, and repeat.  Elementary operations on the active
indices are still lower triangular in the original order.  They do not
change any earlier pivot row or column, whose other entries have already
been cleared.  Each step isolates a rank-one block and leaves the
remaining rank in the active submatrix.  Thus the final partial
permutation matrix has exactly $a$ ones.  If $LAR=\Pi$ is the resulting
identity, then $E_1=L^{-1}$ and $E_2=R^{-1}$ have the required form.
\end{proof}

\begin{lemma}[The cost of a rational matrix shear]
\label{lem:matrix-shear}
Let $\mathcal A$ be a fixed finite collection of rational $m\times m$
matrices.  There is a fixed odd prime $q$ such that the following holds.
Let $H=(H_1,\ldots,H_m)$ and $D=(D_1,\ldots,D_m)$ be two ordered
groups of fields, every field ranging over $[q^b]$, with every $H_i$
before every $D_j$.  For $A\in\mathcal A$ of rational rank $a$, the
permutation
\[
 (H,D)\longmapsto(H+AD,D)
 \quad\text{over }\mathbb Z/q^b\mathbb Z
\]
can be performed with exactly $a$ interchanges of pairs of $b$-digit
fields and $O(V)$ further work.  Every interchange is between an $H$
field and a $D$ field.  Spectators between any of the formal fields are
allowed.
\end{lemma}

\begin{proof}
Choose the factorizations of Lemma~\ref{lem:lower-lower} for all matrices
in $\mathcal A$ once and for all.  Include their factors and inverses in
a finite rational table.  Choose an odd prime avoiding every denominator
in the table and every numerator of a diagonal entry of an invertible
factor.  For example, a prime exceeding the absolute values of all its
nonzero numerators and denominators suffices.  Rational elimination and
trial division construct this finite table and a suitable prime in a
finite computation independent of $b$.  They may therefore be compiled
into one fixed machine description.  All these
rational identities then specialize to $\mathbb Z/q^b\mathbb Z$, and
the triangular factors remain invertible there.

A lower triangular transformation of $H$, or of $D$, costs $O(V)$:
update coordinates in descending physical order.  At coordinate $i$,
multiply by its diagonal unit and add a fixed linear combination of
earlier coordinates.  These earlier coordinates still have their original
values.  Apply Lemma~\ref{lem:ordered-affine-streams}; the number of
coordinates is the fixed number $m$.  Inverses are also lower triangular.

For $A=E_1\Pi E_2$, first transform $D$ by $E_2$ and $H$ by $E_1^{-1}$,
then add $\Pi D$ to $H$, and finally transform $H$ by $E_1$ and $D$ by
$E_2^{-1}$.  Each one in position $(i,j)$ of $\Pi$ calls for
$H_i\leftarrow H_i+D_j$.  Implement it by
\[
 D_j\leftarrow D_j+H_i,\qquad
 (H_i,D_j)\leftarrow(D_j,H_i),\qquad
 D_j\leftarrow H_i-D_j.
\]
The first and last updates target a later field controlled by an earlier
field.  On an initial pair $(h,d)$ this sequence yields $(h+d,d)$.
It uses just one interchange.  There are $a$ ones, which proves the
claim.  In particular, the number of recursive calls is the rational
rank $a$; temporary tapes used for the other operations do not multiply
that count.
\end{proof}

\subsection{Transferring a finite circuit to address permutations}

Fix integers $m,W\ge2$.  The circuit has $W$ wires, including its
scratch wires, with separate input and output terminals.  Its fixed,
pointwise bit gates induce a permutation $\rho$ of the $W$ scalar inputs
at completion: the input value at role $w$ appears at output role
$\rho(w)$.  Assign a
rational $m\times m$ matrix to each terminal and to each gate, the latter
assignment being common to every wire at that gate.  A wire segment
between consecutive assigned vertices is called an edge.  For an edge
$e$ put
\[
 A_e=M_{\mathrm{head}(e)}-M_{\mathrm{tail}(e)},
 \qquad s=\sum_e\operatorname{rank}_{\mathbb Q}A_e.
\]
The required conditions are
\begin{equation}
\label{eq:finite-shear-contract}
 M_{\mathrm{out}(\rho(w))}-M_{\mathrm{in}(w)}=I_m
 \quad(1\le w\le W),\qquad s<Wm.
\end{equation}
All these data are finite and fixed. Proposition~\ref{prop:bit-motif-interface}
supplies these conditions with $W=W_{\rm b}$ and $s=s_{\rm b}$.
Its negative source projection is essential to the endpoint identity and
has already been included in the rank sum.

\begin{proposition}[Power-width interchange with rows]
\label{prop:power-interchange}
Assume the fixed circuit data satisfy
\eqref{eq:finite-shear-contract}, and let $q$ be a prime as in
Lemma~\ref{lem:matrix-shear} for its edge matrices, also avoiding the
denominators of every assigned terminal and gate matrix $M$.
Choose a rational number $\tau\in(0,1)$ with $s/W<m^\tau$.
For every $e=m^k$, $k\ge0$, consider two disjoint contiguous address
chunks of $e$ radix-$q$ digits each.  Suppose a row field preceding both
chunks has length divisible by $W^k$.  There may be additional fields
before and after the row field, between the chunks, and after them.
The chunks can be interchanged in $O(Ve^\tau)$ time, preserving every
other address field, on a fixed number of finite-alphabet tapes.
The constant is independent of all field lengths and of the data.
\end{proposition}

\begin{proof}
For $k=0$ apply the single-digit operation of
Lemma~\ref{lem:elementary-streams}.  For $k>0$, write each chunk as $m$
consecutive $b=e/m$ digit fields, giving vectors $H,D$ over
$\mathbb Z/q^b\mathbb Z$.  We first construct the shear $H\leftarrow H+D$.

Split whole rows cyclically into $W$ role streams within each prefix
preceding the row field.  Write the original row number as
$r=Wg+(w-1)$ with $1\le w\le W$; stream $w$ has row number $g$.
These streams have exactly the same remaining shape, including the
full ranges of $H$ and $D$.  Each has
logical volume $V/W$, and its row count is divisible by $W^{k-1}$.

For a rational matrix $M$, let $\Phi_M$ act on an array by moving its
entry at address $(H,D)$ to $(H+MD,D)$, with every spectator unchanged.
Thus $\Phi_M$ is the array permutation induced by that address map.
Its inverse is $\Phi_{-M}$ and
$\Phi_{M'}\Phi_M=\Phi_{M'+M}$.  On every circuit edge apply
$\Phi_{A_e}$ to that role stream; implement it by
Lemma~\ref{lem:matrix-shear}, recursively interchanging each pivot pair
of $b$-digit fields.  Execute each circuit gate pointwise at matching
addresses of its role streams.

The row split supplies the common address set used in the
frame identity~\eqref{eq:common-frame-identity}. Here the frame at a
vertex with matrix $M$ is the address permutation $\Phi_M$, and every
edge changes it by $\Phi_{M'-M}$. The same invariant proof applies
because a common address permutation commutes with every pointwise
bit gate. For the physical input arrays $f_w$, the logical arrays are
$g_w=\Phi_{-M_{\mathrm{in}(w)}}f_w$; this only specifies their
interpretation, with no initial tape operation. The logical circuit
routes $g_w$ to role $\rho(w)$. Thus its physical output satisfies
\[
 f'_{\rho(w)}
 =\Phi_{M_{\mathrm{out}(\rho(w))}}
   \Phi_{-M_{\mathrm{in}(w)}}f_w
 =\Phi_{I_m}f_w,
\]
by~\eqref{eq:finite-shear-contract}. Merge by returning row $g$ of
output role $\rho(w)$ to the original row $Wg+(w-1)$ within the same
preceding prefix. This undoes $\rho$ and restores row order. The result
is $H\leftarrow H+D$ on every original row, including rows assigned
to scratch roles, whose initial values were arbitrary.

To interchange the two full chunks, perform
\[
 D\leftarrow D-H,\qquad H\leftarrow H+D,\qquad D\leftarrow H-D,
\]
componentwise.  Only the middle shear uses the circuit.  The other two
updates target the later group and have linear cost.  The transformation
sends $(H,D)$ to $(D,H)$.

The number of recursive calls is the sum of the edge ranks, exactly $s$.
Each acts on one role stream of volume $V/W$, retaining all other
fields as spectators.  Splitting,
merging, gates, and triangular operations have total cost $O(V)$ because
their number is fixed.  Below we verify that scheduling the recursive
calls on fixed tapes also costs $O(V)$ at this node.  Thus if $F_k$ is a
uniform upper bound for time divided by logical volume, then
\begin{equation}
\label{eq:swap-recurrence}
 F_0=O(1),\qquad F_k\le(s/W)F_{k-1}+O(1).
\end{equation}
For $a=s/W<m^\tau$ the geometric sum
$\sum_{j=0}^k a^j$ is $O(m^{k\tau})$.  If $a\le1$, the sum is at
most $k+1=O(m^{k\tau})$; if $a>1$, it is $O(a^k)$.
Hence $F_k=O(e^\tau)$.

\paragraph{A fixed-tape depth-first schedule.}
Here we complete the implementation assertion used in
\eqref{eq:swap-recurrence}.  Reserve $W$ role tapes, a fixed set of
I/O and elementary-operation work tapes, one parking stack tape, and a
separate descriptor stack.  Although $W$ is large, it is fixed.
At a call boundary all reusable work tapes, other than the current I/O
and local descriptor, are empty and their heads are at their origins.
The parking and descriptor heads are at the tops of their stacks.

Before a recursive call, append the $W-1$ inactive role streams to the
parking stack in a fixed order, with separators.  Copy the active stream
to the child I/O area, and erase the parent role tapes while traversing
their current contents.  Push the parent descriptor and its finite
program counter to the descriptor stack.  The role and work tapes are
now available to the child, using the same tapes at every depth.
On return, copy the child output into the designated active role.
Pop the inactive streams in reverse order.  To recover a stream in its
original direction, prepare its destination to the known length and
write the popped symbols from the destination's right end toward its
left end.  Erase the popped stack cells as they are left.  This uses
linear time and ends with the parking head at the previously saved top.
No scan passes through a parked ancestor.  Cleanup is always limited to
the region visited by the current call; a reused tape is not swept to
the largest position ever visited by an ancestor.

Each push, pop, copy, positioning operation, or erasure moves $O(V)$
symbols at a parent node.  There are only the fixed number $s$ of child
calls.  This gives precisely the additional $O(V)$ term claimed above,
even though each child itself has volume $V/W$.  Temporary storage is
therefore not substituted for logical volume in the factor $s/W$.

For completeness, descriptor processing also fits this term when the
payload is one bit.  Let $V_0$ be the volume at the root of one
power-width call with width $e_0=m^{k_0}$, and $V_j=V_0/W^j$ the volume
of any depth-$j$ child.  Here $j\le k_0=\log_m e_0$.  The root row
count is at least $W^{k_0}$, so at depth $j$ it is still at least one.
Both original chunk ranges remain present: digits outside the current
subchunks become spectators rather than being removed.  Hence
\[
 V_j\ge q^{2e_0},\qquad
 j\log W\le(\log_m e_0)\log W=O(e_0)=O(\log V_j).
\]
It follows that
$\log V_0=\log V_j+j\log W=O(\log(2V_j))$.
A local descriptor lists the row field, at most the current $m+m$
fields, and a fixed number of intervals formed from intervening
spectator fields.  Computing products of their lengths, subdividing
target fields, copying lengths, and storing the current instruction take
$(\log(2V_0))^{O(1)}=O(V_j)$ time at that child, since every fixed
power of $\log(2V_j)$ is $O(V_j)$.  Before descending,
construct the child's descriptor by retaining its two target fields and
combining each consecutive group of spectator fields into one interval.
The ancestor's subdivision need not remain active in the child.  Stack entries contain
finite instructions and binary integers, not new alphabet symbols.
These observations establish one finite machine for every recursion
depth and every choice of spectator lengths.

\end{proof}


We use the rational value $\tau=1-2^{-50}$ established in
\eqref{eq:explicit-motif-exponents}. Its verification uses fixed integer
comparisons and the displayed exponential inequality; it needs no
real-arithmetic oracle.


\subsection{Removing width and row restrictions}

We have proved the fast operation for power widths with enough rows.
Rows for a general array can be supplied by a few of its own high digits.
Their movement will cost less than the final bound.

\begin{lemma}[Arbitrary-width interchange]
\label{lem:chunk-swap}
There exist fixed $0<\tau<1$ and one fixed finite-alphabet multitape
procedure with the following property.  For positive integers $P,G,B$
and $u\ge1$, it transforms an arbitrary bit array on
\[
 [P]\times[2^u]\times[G]\times[2^u]\times[B]
\]
by the address permutation
\[
 (p,h,g,d,z)\longmapsto(p,d,g,h,z)
\]
in $O(Vu^\tau)$ steps, where $V=PG B\,2^{2u}$.
The integers specifying the shape are part of the input.  The time
includes their processing, all padding and unpadding, and fixed-tape
workspace cleanup.  In particular, $B=1$ is allowed.
\end{lemma}

\begin{proof}
First suppose the two ranges are $[q^e]$, where $q$ is the fixed prime
already chosen.  Set
\[
 k=\lceil\log_m e\rceil,\qquad
 \rho=\left\lceil\frac{k\log W}{2\log q}\right\rceil.
\]
For all sufficiently large $e$, $1\le\rho<e$ and $\rho=O(\log(2e))$.
Handle the bounded remaining set of widths by $e$ individual digit
interchanges; the constant in $O(Ve^\tau)$ absorbs their cost.

Write each chunk as its top $\rho$ digits followed by its remaining
digits, denoted $h^+,h^-$ and $d^+,d^-$, respectively.  Interchange
the corresponding high digits one pair at a time.  The high positions
now contain $d^+$ and $h^+$; move the digits of the latter, preserving
their order, into consecutive positions immediately after $d^+$.
Suppressing the prefix and suffix, the field orders during the construction are
\[
\begin{aligned}
 &(h^+,h^-,g,d^+,d^-),\\
 &(d^+,h^+,h^-,g,d^-)
       &&\text{after exchanging and joining the high parts},\\
 &(d^+,h^+,d^-,g,h^-)
       &&\text{after interchanging the remaining parts},\\
 &(d^+,d^-,g,h^+,h^-)
       &&\text{after separating the high parts again}.
\end{aligned}
\]
In the two middle orders, the first two fields form a row field of range
\[
 R=q^{2\rho}\ge W^k
\]
before both remaining chunks.  The joining moves preserve the order of
the remaining digits and the gap field.  Exchanging and joining the high
parts costs $O(V\rho)$ by
Lemma~\ref{lem:elementary-streams}.  Append zero rows, separately within
each preceding prefix, until the row count is
\[
 R'=W^k\left\lceil R/W^k\right\rceil<2R.
\]
This changes the row field's length; it need not remain a power of $q$
during the computation.  The padded logical volume is less than $2V$.

Write the remaining width in base $m$:
\[
 e-\rho=\sum_{j=0}^{k} a_jm^j,\qquad 0\le a_j<m.
\]
Partition each remaining chunk into corresponding consecutive pieces,
with $a_j$ pieces of width $m^j$.  Apply
Proposition~\ref{prop:power-interchange} to each corresponding pair.
All these calls use the same row field of length $R'$, which is divisible
by each required $W^j$.  At completion every call preserves its row
address, so it neither mixes genuine rows with padded rows nor depends
on the values used for padding.  The padded zero rows remain zero at
these completed-call boundaries.  The sum of the costs is bounded by
\[
 O(V)\sum_j a_jm^{j\tau}=O(Ve^\tau).
\]
The fixed bound $a_j<m$ and the inequality $m^k<me$ make this a
geometric sum of the stated order.  Scan away the padded rows, then
reverse only the moves that joined the high parts.  This is the last
transition in the field-order display: $h^+$ returns to the later
high position, while $d^+$ stays in the earlier one.  The final order
is therefore the required interchange of the two full chunks.  The
remaining movement cost $O(V\log(2e))$ is $O(Ve^\tau)$.
Computing the piece list and every new shape is polynomial in
$\log(2V)$ per piece.  There are at most
$(m-1)(k+1)=O(\log(2e))$ pieces; since $V\ge q^{2e}$, their total
descriptor cost is a fixed power of $\log(2V)$ and hence $O(V)$.
Each power-width call uses the fixed schedule proved above and leaves
its work tapes clean.

Finally, for binary chunks let $e$ be the least integer with
$q^e\ge2^u$.  Pad each of the two chunk ranges numerically from $[2^u]$
to $[q^e]$, putting zero in an entry if either chunk coordinate is
invalid.  The increase in volume is less than $q^2$, a fixed factor.
The padding is a nested scan through consecutive valid and invalid
intervals, with their lengths prepared from the shape.  It costs $O(V)$
with a constant depending on $q$.  Lexicographic radix-$q$ order of
$[q^e]$ is its numerical order, just as binary order of $[2^u]$ is its
numerical order.  Padding therefore preserves the order of the existing
entries without converting the digits of every address.
The completed interchange preserves the set on which both chunk
coordinates lie below $2^u$.  Scan away the other entries to recover the
desired binary array.  Since $e=O(u)$ for fixed $q$, the bound is
$O(Vu^\tau)$.
\end{proof}

The lemma permits arbitrary prefix, gap, and suffix lengths.  Thus every
later interchange of equal bit chunks is an instance of its stated
layout, after multiplying the irrelevant contiguous lengths.
For widths differing by one bit, choose the leading bit of the longer
chunk as the extra bit.  Suppress the outer prefix and suffix, and write
$g$ for the intervening address field.  If the first chunk is longer,
write it as $(h_0,H)$, so $H$ and the second chunk $D$ have equal width.
The physical orders are
\[
 (h_0,H,g,D)\longmapsto(h_0,D,g,H)
              \longmapsto(D,g,h_0,H).
\]
The first step interchanges $H,D$ by the lemma; the second moves the
single bit $h_0$ immediately before $H$.  If the second chunk is longer,
write it as $(d_0,D)$, with $H$ the first chunk.  Use
\[
 (H,g,d_0,D)\longmapsto(d_0,H,g,D)
              \longmapsto(d_0,D,g,H).
\]
Here the first step moves $d_0$ immediately before $H$, and the second
interchanges $H,D$.  In each case $g$ remains between the two full chunks
in the final order.  If $H,D$ are empty, the indicated equal-width
interchange is the identity.  Each case uses one single-bit move, costing
$O(V)$ by Lemma~\ref{lem:elementary-streams}, in addition to the
equal-width interchange.  No unrestricted array transpose is being used
in these operations.
